% !TeX root = ../../Book5.tex
% 纲要标题：### 11.3 表示与模
% 纲要位置：工作记录/计划纲要.md，第 4099 行。

\section{表示与模}
\label{b5:sec:1103}


% 纲要范围：
%
% * Lie 代数的线性作用
% * 对偶和张量表示
% * 不变子空间及完全可约性的提问
%

Lie 代数的表示仍然由线性算子组成，但相容条件针对的是括号。一个元素无需作用为可逆算子，也没有“先求群元素的逆”的步骤。对偶与张量作用中的符号应从这个条件重新推导，不能照搬群表示的公式。

\subsection{Lie 代数的线性作用}
\label{b5:subsec:1103:01}

\begin{definition}\label{b5:def:lie-representation}
设 \(\mathfrak g\) 为域 \(k\) 上的 Lie 代数，\(V\) 为 \(k\) 向量空间。一个 \(k\) 线性映射
\(\rho:\mathfrak g\to\mathfrak{gl}(V)\)
若满足
\[
\rho\bigl([x,y]\bigr)=\rho(x)\rho(y)-\rho(y)\rho(x)
\quad(x,y\in\mathfrak g),
\]
就称为 \(\mathfrak g\) 在 \(V\) 上的\term[lie-daishu-biaoshi]{表示}[representation of a Lie algebra]，也称 \(V\) 为左 \(\mathfrak g\) 模。简记 \(xv=\rho(x)v\)；此时相容条件为
\[
[x,y]v=x(yv)-y(xv).
\]
\end{definition}
设 \(\rho_V,\rho_W\) 分别为 \(\mathfrak g\) 在 \(V,W\) 上的表示。线性映射 \(T:V\to W\) 若满足
\[
T\rho_V(x)=\rho_W(x)T\quad(x\in\mathfrak g),
\]
即 \(T(xv)=xT(v)\)，就称为表示同态；两个表示同构要求这样的同态可逆。当 \(V=W\) 且两表示相同时，这一条件才表示 \(T\) 与全部作用算子对易。子表示是被全部作用算子保持的子空间，商表示由商空间上的诱导算子给出。若 \(\rho\) 单射，称表示忠实；若全部作用为零，称表示平凡。
\ProseParagraphBreak
\cref{b5:prop:adjoint-representation}给出了伴随表示。矩阵 Lie 代数作为 \(\mathfrak{gl}_n(k)\) 的子代数时，还在 \(k^n\) 上有自然表示。另一个简单例子是一维模：它由线性泛函 \(\lambda:\mathfrak g\to k\) 决定，作用为 \(xv=\lambda(x)v\)。表示条件恰好要求
\(\lambda\bigl([\mathfrak g,\mathfrak g]\bigr)=0\)，因为一维算子彼此交换。故一维模与交换化商 \(\mathfrak g/[\mathfrak g,\mathfrak g]\) 上的线性泛函对应。
\begin{proposition}\label{b5:prop:lie-invariant-subquotients}
设 \(\mathfrak g\) 为域 \(k\) 上的 Lie 代数，\(V,W\) 为左 \(\mathfrak g\) 模，\(f:V\to W\) 为表示同态。则核、像与余核均自然为左 \(\mathfrak g\) 模；一列左 \(\mathfrak g\) 模正合，当且仅当底层 \(k\) 向量空间列正合。
\end{proposition}
\begin{proof}
若 \(f(v)=0\)，则 \(f(xv)=xf(v)=0\)，故核不变。若 \(w=f(v)\)，则 \(xw=f(xv)\) 仍在像中。于是商上的作用不依赖代表，且表示等式从原空间下降。核、像和商都没有改变相应的向量空间运算，所以正合性也按底层空间检验。
\end{proof}

\subsection{对偶表示与张量表示}
\label{b5:subsec:1103:02}

若要求求值 \(V^*\otimes V\to k\)、\(\varphi\otimes v\mapsto\varphi(v)\) 是表示同态，而 \(k\) 上的作用平凡，那么必有
\[
0=x\cdot\varphi(v)=(x\varphi)(v)+\varphi(xv).
\]
这便迫使对偶作用带负号。张量作用的两项之和则可从一次项看出：在 \(\varepsilon^2=0\) 时，
\[
(I+\varepsilon\rho_V(x))v\otimes(I+\varepsilon\rho_W(x))w
=v\otimes w+\varepsilon(xv\otimes w+v\otimes xw).
\]
下面还须核验这些由相容性要求得到的公式确实满足 Lie 括号关系。

\begin{proposition}\label{b5:prop:lie-dual-tensor-hom}
设 \(\mathfrak g\) 为域 \(k\) 上的 Lie 代数，\(\rho_V,\rho_W\) 分别为它在 \(k\) 向量空间 \(V,W\) 上的表示。对 \(x\in\mathfrak g\)、\(v\in V\)、\(w\in W\)、\(\varphi\in V^*=\operatorname{Hom}_k(V,k)\) 及 \(T\in\operatorname{Hom}_k(V,W)\)，下列公式分别给出 \(V^*\)、\(V\otimes_k W\) 与 \(\operatorname{Hom}_k(V,W)\) 上的表示；其中 \(xv=\rho_V(x)v\)、\(xw=\rho_W(x)w\)：
\[
\begin{aligned}
(x\varphi)(v)&=-\varphi(xv),\\
x(v\otimes w)&=xv\otimes w+v\otimes xw,\\
xT&=\rho_W(x)T-T\rho_V(x).
\end{aligned}
\]
若 \(V\) 有限维，则自然同构 \(V^*\otimes W\cong\operatorname{Hom}_k(V,W)\) 与这些作用相容。
\end{proposition}
\begin{proof}
在对偶上，先后作用 \(x,y\) 给出
\[
\begin{aligned}
\bigl(x\cdot(y\cdot\varphi)\bigr)(v)&=\varphi\bigl(y(xv)\bigr),\\
\bigl(y\cdot(x\cdot\varphi)\bigr)(v)&=\varphi\bigl(x(yv)\bigr).
\end{aligned}
\]
由 \(V\) 上的表示关系，
\[
\begin{aligned}
\bigl(x\cdot(y\cdot\varphi)-y\cdot(x\cdot\varphi)\bigr)(v)
&=-\varphi\bigl([x,y]v\bigr)\\
&=\bigl([x,y]\cdot\varphi\bigr)(v),
\end{aligned}
\]
正好是对偶上的表示关系。在张量积上，算子为
\(\rho_V(x)\otimes I+I\otimes\rho_W(x)\)；
两个不同因子上的算子交换，计算交换子只留下两个因子的表示关系。对 Hom 展开两次作用，混合项相消，同样得到
\(\rho_W\bigl([x,y]\bigr)T-T\rho_V\bigl([x,y]\bigr)\)。
最后把 \(\varphi\otimes w\) 送到 \(v\mapsto\varphi(v)w\)，代入公式得到
\[
x\bigl(\varphi(v)w\bigr)-\varphi(xv)w
=\varphi(v)xw-\varphi(xv)w,
\]
这正是对偶与张量作用的像。有限维线性代数给出该映射为同构。
\end{proof}
尤其
\[
\operatorname{Hom}_{\mathfrak g}(V,W)
=\operatorname{Hom}_k(V,W)^{\mathfrak g},
\qquad
U^{\mathfrak g}=\{\,u\in U\mid xu=0\text{ 对所有 }x\in\mathfrak g\,\}.
\]
这里“不变”表示被所有 Lie 代数作用零化，与群表示中被所有群元素固定的说法不同。
\ProseParagraphBreak
在 \(V^{\otimes m}\) 上，\(x\) 的作用是分别在各因子作用后相加，并令 \(x1=0\)。将两个张量字拼接时，这些因子分为两组，因而在张量代数 \(T(V)\) 上有
\[
x(u\otimes w)=(xu)\otimes w+u\otimes(xw)
\quad(u,w\in T(V)).
\]
因此每个 \(x\) 在 \(T(V)\) 上作用为导子。对称代数的关系是由 \(v\otimes w-w\otimes v\) 生成的双边理想，而
\[
\begin{aligned}
x(v\otimes w-w\otimes v)
&=(xv\otimes w-w\otimes xv)\\
&\quad +(v\otimes xw-xw\otimes v)
\end{aligned}
\]
仍属于该理想。由导子性质，这个双边理想被作用保持。

外代数的关系理想由平方张量生成。对 \(v\otimes v\) 的作用为
\[
\begin{aligned}
x(v\otimes v)&=xv\otimes v+v\otimes xv\\
&=(xv+v)\otimes(xv+v)-xv\otimes xv-v\otimes v,
\end{aligned}
\]
所以该理想也被保持。因此对称幂、外幂及整个对称代数和外代数都得到诱导作用。在多项式代数上，这个作用满足 Leibniz 规则\footnote{莱布尼茨（Gottfried Wilhelm Leibniz，1646-1716），德国数学家、哲学家，独立创立微积分，并发展了微分记号和乘积求导法则。}，是导子，而不是代数自同构。

\subsection{不变子空间与完全可约性}
\label{b5:subsec:1103:03}

\begin{definition}\label{b5:def:lie-simple-semisimple-module}
设 \(\mathfrak g\) 为域 \(k\) 上的 Lie 代数，\(V\) 为左 \(\mathfrak g\) 模。若 \(V\ne0\) 且没有非零真子表示，就称 \(V\) 为不可约表示；若 \(V\) 是不可约子表示的直和，就称为完全可约表示。若 \(V\ne0\) 且不能分成两个非零子表示的直和，则称为不可分解表示。有限维不可约表示必不可分解，逆命题一般不成立。
\end{definition}
例如二维可解 Lie 代数 \([h,e]=e\) 可在 \(k^2\) 上作用为
\[
\rho(h)=\begin{pmatrix}1&0\\0&0\end{pmatrix},\qquad
\rho(e)=\begin{pmatrix}0&1\\0&0\end{pmatrix}.
\]
两矩阵的交换子等于 \(\rho(e)\)，所以确为表示。第一坐标轴不变，但任何补直线均含第二坐标非零的向量，而 \(\rho(e)\) 将它送到非零第一坐标向量，因而没有不变补。这个表示可约却不可分解。
\ProseParagraphBreak
Schur 引理对 Lie 代数表示仍成立：同态的核与像是子表示，所以不可约之间的非零同态可逆。若基域代数闭且表示有限维，自同态取一个特征值后减去相应标量，就由同一论证得到它原本就是标量。但要从自同态只有标量反推不可约，仍须有额外条件；上一例的两个作用矩阵的公共中心化子就只有标量，却不完全可约。
\ProseParagraphBreak
Lie 代数表示的分解取决于 Lie 代数本身的结构。本编将证明：半单复 Lie 代数的有限维表示完全可约，可解复 Lie 代数的有限维不可约表示只能一维。前一结论需要 Killing 型与 Casimir 方法，后一结论来自同时三角化；两者的假设与证明均不同。
